\section{Tổng quan}

Repo \repo{} hiện thực bài toán sinh ảnh bằng \textit{Shortcut Models} trên nền JAX/Flax.
Ý tưởng chính là huấn luyện một mô hình DiT duy nhất để vừa học mục tiêu flow matching tiêu chuẩn,
vừa học thêm các shortcut target ở nhiều độ dài bước khác nhau. Nhờ đó, cùng một checkpoint có thể
được dùng cho nhiều ngân sách lấy mẫu, từ nhiều bước xuống một bước.

\subsection{Mục tiêu của repo}

Ba mục tiêu vận hành chính của repo là:

\begin{itemize}
    \item huấn luyện Shortcut Models theo bài báo gốc;
    \item giữ các baseline để so sánh, gồm \code{naive}, \code{sit}, \code{progressive}, \code{consistency},
    \code{consistency-distillation} và \code{livereflow};
    \item cung cấp cùng một entrypoint cho train, đánh giá định kỳ và sampling offline.
\end{itemize}

\subsection{Các mode huấn luyện}

Repo hiện hỗ trợ các giá trị \code{--model.train_type} sau:

\begin{itemize}
    \item \code{shortcut}: mode chính của bài báo. Batch được trộn giữa flow target và bootstrap shortcut target.
    \item \code{naive}: flow matching cơ bản, không dùng shortcut target.
    \item \code{sit}: baseline Scalable Interpolant Transformers dùng cùng backbone DiT nhưng đổi sang transport continuous-time với path, prediction target và sampler cấu hình được.
    \item \code{progressive}: dùng teacher state để rút ngắn dần bước theo tiến độ train.
    \item \code{consistency}: consistency training dựa trên EMA.
    \item \code{consistency-distillation}: teacher-student consistency distillation.
    \item \code{livereflow}: sinh một phần target kiểu reflow trực tiếp từ rollout của model hiện tại.
\end{itemize}

\begin{figure}[H]
    \centering
    \includegraphics[width=0.92\textwidth]{fig-method5.png}
    \caption{Hình minh hoạ trực quan cho Shortcut Models trong tài liệu gốc.}
\end{figure}

\subsection{Phạm vi tài liệu PDF này}

Tài liệu tập trung vào các nội dung cần cho việc vận hành repo:

\begin{itemize}
    \item file nào là entrypoint chính và file nào chứa target/model;
    \item luồng dữ liệu từ dataset sang latent, target và update;
    \item lệnh thiết lập môi trường, train và inference phù hợp với code hiện tại;
    \item các điểm dễ gây nhầm lẫn giữa tài liệu cũ và hành vi thực tế của nguồn.
\end{itemize}
